\documentclass[10pt,a4paper]{amsart}
\usepackage[margin=19mm]{geometry}
\usepackage{amsmath,amssymb,mathtools}
\usepackage[scr=rsfs,cal=euler]{mathalfa}
\usepackage{xfrac}
\usepackage[unicode,hidelinks,bookmarks=false]{hyperref}
\newcommand{\ie}{i.e.\ }
\newcommand{\cf}{cf.\ }
\newcommand{\resp}{respectively\ }
\newcommand{\Subsection}{\subsection}
\providecommand{\nobreakdash}{\nobreak\hbox{-}}
\setlength{\emergencystretch}{2em}
\multlinegap0pt
\usepackage[shortlabels]{enumitem}
\usepackage{amssymb}
\newtheorem{prop}{Proposition}
\def\mod#1{{\ifmmode\text{\rm\ (mod~$#1$)}
\else\discretionary{}{}{\hbox{ }}\rm(mod~$#1$)\fi}}
\title{Perfect $2$-codes over arbitrary alphabets}
\author{Michael A. Bennett}
\date{Source: arXiv:2607.27555v1 (CC BY 4.0)}
\begin{document}
\maketitle

\noindent\emph{Source and licence.} Perfect $2$-codes over arbitrary alphabets. Version arXiv:2607.27555v1 (CC BY 4.0), distributed by arXiv under the Creative Commons Attribution 4.0 International licence, http://creativecommons.org/licenses/by/4.0/, as recorded in the arXiv licence metadata for this exact version and shown on the abstract page https://arxiv.org/abs/2607.27555v1. Full licence notice: \texttt{licenses/arxiv-2607.27555-CC-BY-4.0.txt}.\par\smallskip

\noindent\textit{Editorial scope.} Counted unit: the article's prop Goop (source0.tex lines 284--289). The article's complete original proof of it is delivered with it (source0.tex lines 291--325, one \texttt{proof} environment of 35 lines). No mathematics is written, repaired or re-derived: the statement and the proof are the article's own lines, in the article's own notation, and no retained line is substituted or re-flowed.  Retained uncounted as the source-local context the proof needs: lines 236--238; lines 240--242; lines 266--268.  Nothing was omitted from any retained span, so the package carries no omission record.  No retained line contains a citation, so the wrapper carries no bibliography and no bibliography entry.  No retained line contains a figure environment, an image-inclusion command or drawing code, so no image is delivered and the package declares no asset dependency.  Editorial formatting only, in addition: an AMS \texttt{amsart} preamble that restates the article's own declarations and macros used by the retained lines, namely the environments \texttt{prop} on separate counters, together with the standard packages those lines need.  The retained environments print in this document as Proposition 1 in order of appearance, whereas the article prints the same items with the numbering of its own full document.  The complete native source of the article as distributed in the arXiv e-print for this exact version is bundled unchanged as \texttt{sources/206352/source0.tex}.
\begin{equation} \label{poly}
x^2 - \left( \left( 2 -\frac{2}{q} \right) (n-2) +3 \right)x+\frac{2f(n,q)}{q^2}
\end{equation}

\begin{equation} \label{roots1}
r_1+r_2=\left( 2 -\frac{2}{q} \right) (n-2) +3
\end{equation}

\begin{equation} \label{good}
q (r_2-r_1)^2 = 2(r_1+r_2) +q-6.
\end{equation}

\begin{prop} \label{Goop}
If there exists a perfect $2$-code on $q$ symbols, with words of length $n>2$, then the (integer) roots $r_1<r_2$ of the polynomial (\ref{poly}) satisfy
\begin{equation} \label{size}
0 < r_2-r_1 < \left( \frac{6r_1}{q} \right)^{1/2}.
\end{equation}
\end{prop}

\begin{proof}
Since $n > 2$,  from (\ref{roots1}),  $n \equiv 2 \mod{q/2}$ if $q$ is even, and  $n \equiv 2 \mod{q}$ if $q$ is odd. Further, from (\ref{roots1}),
$$
r_1+r_2 > \left( 2 -\frac{2}{q} \right) n  -1 
$$
and hence, from $r_2 \leq n$, 
\begin{equation} \label{r1low}
r_1 > \left( 1 -\frac{2}{q} \right) n  -1 > 2,
\end{equation}
where we have used only that $q \geq 6$ and $n \geq 5$. Suppose then that (\ref{size}) fails to hold, so that we may write $r_2=r_1+\frac{\kappa}{\sqrt{q}} \,  r_1^{1/2}$, with $\kappa \geq \sqrt{6}$. Then (\ref{good}) implies that
$$
\kappa^2  \, r_1 = 4 r_1+ \frac{2 \kappa}{\sqrt{q}} r_1^{1/2} +q-6,
$$
whence
\begin{equation} \label{what}
\frac{2}{\sqrt{6}} \leq \frac{\kappa^2-4}{\kappa} = \frac{2}{\sqrt{q r_1}} + \frac{q-6}{\kappa r_1}.
\end{equation}
If $q$ is even, $n \equiv 2 \mod{q/2}$ and $n > 2$ imply that  $n \geq 2+q/2$. Thus
$$
r_1 > \left( 1 -\frac{2}{q} \right) n  -1 \geq 
\left( 1 -\frac{2}{q} \right) (2+q/2)-1 =q/2-4/q
$$
and so $r_1 \geq q/2$. If $q$ is odd, $n \equiv 2 \mod{q}$, whence
we have $n \geq 2+q$ and so 
$$
r_1 > \left( 1 -\frac{2}{q} \right) n  -1 \geq 
\left( 1 -\frac{2}{q} \right) (2+q)-1 =q-4/q-1,
$$
and again, crudely, $r_1 \geq q/2$. Thus, from (\ref{what}) and $\kappa \geq 6$, 
$$
\frac{2}{\sqrt{6}}  
< \frac{2\sqrt{2}}{q} + \frac{2}{\sqrt{6}}  -\frac{12}{\sqrt{6} q},
$$
a contradiction.
\end{proof}

\end{document}
